% Printed182; principalPDF189.
\par\clearpage
\noindent182\hfill{\small\textbf{POSTHUMOUS PAPERS.}}\hfill\null\par\medskip
It remains to expand \(\frac{a\cos\varphi}{\sqrt{1-ee\sin\varphi^2}}\), which is obtained from the formula \(-\int\mathrm{d}\psi.\sin\varphi\) [multiplied by the factor \(K=a(1-4ff-12f^4-80f^6-700f^8\ldots)\)].

From this we have found the following value:
\[
\begin{aligned}
\left[\frac{a\cos\varphi}{K\sqrt{1-ee\sin\varphi^2}}={}\right]
&\,(1+3ff+15f^4+87f^6+528f^8\ldots)\cos\psi\\
&+(ff+10f^4+88f^6+745f^8\ldots)\cos3\psi\\
&+\left(3f^4+\frac{158}{3}f^6+672f^8\ldots\right)\cos5\psi\\
&+\left(\frac{37}{3}f^6+312f^8\ldots\right)\cos7\psi\\
&+(59f^8\ldots)\cos9\psi
\end{aligned}
\]
{\centering etc.\par}
\bigskip
{\centering [28.]\par}
[For \(n=302{,}78\) we have]
\[
\log ff=6{,}615\,0650.573
\]
\[
\begin{array}{@{}r@{\;}c@{\;}l@{\qquad}r@{\;}c@{\;}r@{}}
[\text{or }\log ff\dfrac{2.10^7}{\pi}&=&3{,}418\,9451.803;&[\dfrac{2.10^7}{\pi}ff&=&2623{,}887\,3172\\[3pt]
\log f^4\quad\text{»}&=&0{,}034\,0102.376;&\text{»}\quad f^4&=&1{,}081\,4594\\
\log f^6\quad\text{»}&=&6{,}649\,0752.949-10;&\text{»}\quad f^6&=&0{,}000\,4457\\
\log f^8\quad\text{»}&=]&3{,}264\,1403.522-10;&\text{»}\quad f^8&=]&0{,}000\,0002.
\end{array}
\]
[For the following table, however, it was assumed that]
\[
\log\frac{2.10^7}{\pi}ff=3{,}418\,9411.803;\qquad
\frac{2.10^7}{\pi}ff=2623{,}863\,1498,
\]
\noindent [to which \(n=302{,}7827\ldots\) would correspond. For this value of \(ff\), the following relations hold between \(\psi,\varphi\) and \(\omega\) when these quantities are measured in metres, that is, in \(10\,000\,000\)ths of the quadrant:]
\[
\begin{aligned}
\varphi-\psi={}&15795^{\mathrm{m}}{,}278\,6227\sin2\psi\\
&+22{,}860\,7756\sin4\psi\\
&+0{,}045\,3131\sin6\psi\\
&+0{,}000\,1008\sin8\psi
\end{aligned}
\]
{\centering etc.\par}

% Printed183; principalPDF190.
\par\clearpage
\noindent\null\hfill{\small\textbf{CONFORMAL MAPPING OF THE SPHEROID IN THE PLANE.}}\hfill183\par\medskip
\[
\begin{aligned}
\omega-\psi={}&-5253^{\mathrm{m}}{,}500\,2000\sin2\psi\\
&-0{,}364\,7702\sin4\psi\\
&-0{,}001\,0182\sin6\psi\\
&-0{,}000\,0013\sin8\psi
\end{aligned}
\]
{\centering etc.\par}
[The table on p.~184 was calculated using these two equations. The other relations between \(\psi,\varphi\) and \(\omega\) are:]
\[
\begin{aligned}
\psi-\omega={}&+5253^{\mathrm{m}}{,}498\,1263\sin2\omega\\
&+4{,}700\,0528\sin4\omega\\
&+0{,}007\,2875\sin6\omega\\
&+0{,}000\,0145\sin8\omega
\end{aligned}
\]
{\centering etc.\par}
\[
\begin{aligned}
\omega-\varphi={}&-21048^{\mathrm{m}}{,}753\,7173\sin2\varphi\\
&+28{,}998\,8622\sin4\varphi\\
&-0{,}049\,9078\sin6\varphi\\
&+0{,}000\,0923\sin8\varphi
\end{aligned}
\]
{\centering etc.\par}
\[
\begin{aligned}
\varphi-\omega={}&+21048^{\mathrm{m}}{,}734\,5472\sin2\omega\\
&+40{,}595\,1921\sin4\omega\\
&+0{,}107\,3684\sin6\omega\\
&+0{,}000\,3194\sin8\omega
\end{aligned}
\]
{\centering etc.\par}
\[
\begin{aligned}
\psi-\varphi={}&-15795^{\mathrm{m}}{,}286\,7250\sin2\varphi\\
&+16{,}329\,1611\sin4\varphi\\
&-0{,}021\,0061\sin6\varphi\\
&+0{,}000\,0289\sin8\varphi
\end{aligned}
\]
{\centering etc.\par}
The radius of the equator [is then, when \(ff,f^4,\ldots\) are already expressed in metres,]
\[
\frac{20\,000\,000}{\pi}+4ff+28f^4+240f^6+2316f^8+24240f^{10}+\ldots
=637\,6723{,}563\,9821\,\mathrm{m}.
\]

% Printed184; principalPDF191.
\par\clearpage
\noindent184\hfill{\small\textbf{POSTHUMOUS PAPERS.}}\hfill\null\par\medskip
{\small\renewcommand{\arraystretch}{.82}
\[\begin{array}{|r|r|r|r|r|r|r|}\hline
\rule{0pt}{2.1ex}\psi&\varphi&\multicolumn{2}{c|}{\text{diff.}}&\omega&\multicolumn{2}{c|}{\text{diff.}}\\
\rule[-3.6pt]{0pt}{3.6pt}\text{(metres)}&\text{(metres)}&\multicolumn{2}{c|}{}&\text{(metres)}&\multicolumn{2}{c|}{}\\\hline\hline
\rule{0pt}{2.5ex}578\,0000&579\,5312{,}62571&&&577\,4903{,}61218&&\\
&&9987{,}75983&&&10004{,}03110&\\
579\,0000&580\,5300{,}38554&&15068&578\,4907{,}64328&&5026\\
\cline{1-2}\cline{5-5}
&&60915&&&08136&\\
580\,0000&581\,5287{,}99469&&56&579\,4911{,}72464&&21\\
&&45859&&&13157&\\
581\,0000&582\,5275{,}45328&&42&580\,4915{,}85621&&17\\
&&30817&&&18174&\\
582\,0000&583\,5262{,}76145&&31&581\,4920{,}03795&&13\\
&&15786&&&23187&\\
583\,0000&584\,5249{,}91931&&17&582\,4924{,}26982&&10\\
&&9987{,}00769&&&28197&\\
584\,0000&585\,5236{,}92700&&15004&583\,4928{,}55179&&04\\
\cline{1-2}\cline{5-5}
&&9986{,}85765&&&33201&\\
585\,0000&586\,5223{,}78465&&14990&584\,4932{,}88380&&5001\\
&&70775&&&38202&\\
586\,0000&587\,5210{,}49240&&78&585\,4937{,}26582&&4996\\
&&55797&&&43198&\\
587\,0000&588\,5197{,}05037&&63&586\,4941{,}69780&&91\\
&&40834&&&48189&\\
588\,0000&589\,5183{,}45871&&50&587\,4946{,}17969&&88\\
&&25884&&&53177&\\
589\,0000&590\,5169{,}71755&&35&588\,4950{,}71146&&83\\
\cline{1-2}\cline{5-5}
&&9986{,}10949&&&58160&\\
590\,0000&591\,5155{,}82704&&22&589\,4955{,}29306&&78\\
&&9985{,}96027&&&63138&\\
591\,0000&592\,5141{,}78731&&14908&590\,4959{,}92444&&74\\
&&81119&&&68112&\\
592\,0000&593\,5127{,}59850&&14893&591\,4964{,}60556&&69\\
&&66226&&&73081&\\
593\,0000&594\,5113{,}26076&&78&592\,4969{,}33637&&64\\
&&51348&&&78045&\\
594\,0000&595\,5098{,}77424&&64&593\,4974{,}11682&&60\\
\cline{1-2}\cline{5-5}
&&36484&&&83005&\\
595\,0000&596\,5084{,}13908&&50&594\,4978{,}94687&&55\\
&&21634&&&87960&\\
596\,0000&597\,5069{,}35542&&33&595\,4983{,}82647&&50\\
&&9985{,}06801&&&92910&\\
597\,0000&598\,5054{,}42343&&20&596\,4988{,}75557&&46\\
&&9984{,}91981&&&10004{,}97856&\\
598\,0000&599\,5039{,}34324&&14804&597\,4993{,}73413&&40\\
&&77177&&&10005{,}02796&\\
599\,0000&600\,5024{,}11501&&14788&598\,4998{,}76209&&4935\\
\cline{1-2}\cline{5-5}
&&9984{,}62389&&&10005{,}07731&\\
600\,0000&601\,5008{,}73890&&&599\,5003{,}83940&&\\
\end{array}\]
}
\begin{center}[The table is based on the flattening \(\frac1{302{,}7827\ldots}\).]\end{center}


% Printed185; principalPDF192.
\par\clearpage
\noindent\null\hfill{\small\textbf{CONFORMAL MAPPING OF THE SPHEROID IN THE PLANE.}}\hfill185\par\medskip
{\centering [29.]\par\medskip
Calculation of the [plane rectangular] coordinates [from the geographical\\
coordinates using the series relating \(\varphi,\psi\) and \(\omega\)].\par\smallskip
Test calculation for Varel.\par}
\[
\left.\begin{aligned}
\lambda&=1^\circ\,48'\,24''{,}7109\\
\Phi&=53^\circ\,23'\,57''{,}0322
\end{aligned}\right\}\quad
\begin{gathered}
\text{[Coordinates of Varel. For Göttingen,}\\
\lambda=0\text{ is assumed.]}
\end{gathered}
\]
\[
\left[=\frac{53{,}399\,175\,611}{90}.10^7\right]=5933241{,}735\,\mathrm{m}.
\]
From this [by looking up the value of \(\Phi\) in the column for \(\varphi\) in the preceding table and reading off the corresponding \(\omega\):]
\[
\Omega=5913075{,}164\,\mathrm{m}=53^\circ\,13'\,3''{,}6353.
\]
[\(\Omega\) can also be calculated from the series for \(\omega-\varphi\), p.~183.

Now let \(x',y'\) be the rectangular coordinates of the point in the plane corresponding to the point \(\Omega,\lambda\) on the sphere with the same meridian circumference as the spheroid, when the sphere is mapped conformally onto the plane in such a way that the meridian \(\lambda=0\) becomes the principal meridian. Thus \(\mathrm{f}(x'+iy')=\int\frac{\mathrm{d}\Omega}{\cos\Omega}+i\lambda=\log\operatorname{hyp}\operatorname{tang}(45^\circ+\frac12\Omega)+i\lambda\). To calculate \(x',y'\), one can use the formulas for \textsc{Mercator}'s projection:
\[
\begin{gathered}
\operatorname{tang}c=\operatorname{tang}\lambda\sin\Omega\\
\sin Q=\sin\lambda\cos\Omega\\
\operatorname{tang}Q=\frac{\sin c}{\operatorname{tang}\Omega}\,]\\
\operatorname{tang}\tfrac12(x'-\Omega)=\operatorname{tang}\tfrac12c\operatorname{tang}\tfrac12Q\\
{[}x'-\Omega=\tfrac12(Q+y')\operatorname{tang}\tfrac12c\ldots\\
y'=Q+\frac{Q^3}{6rr}+\frac{Q^5}{24r^4}+\ldots,
\end{gathered}
\]
\noindent where \(r=\frac{20\,000\,000}{\pi}\)].
\[
\begin{array}{@{}l@{\;}r@{\quad}l@{\;}r@{\quad}l@{\;}r@{}}
\sin\Omega\,.\,.\,.\,.&9{,}903\,5870&\cos\Omega\,.\,.\,.\,.&9{,}777\,2648&\sin c\,.\,.\,.\,.\,.&8{,}402\,3953\\
\operatorname{tang}\lambda\,.\,.\,.&8{,}498\,9469&\sin\lambda\,.\,.\,.\,.&8{,}498\,7309&\operatorname{tang}\Omega\,.\,.&0{,}126\,3223\\\cline{2-2}\cline{4-4}\cline{6-6}
\operatorname{tang}c\,.\,.\,.&8{,}402\,5339&\sin Q\,.\,.\,.\,.&8{,}275\,9957&\operatorname{tang}Q\,.\,.&8{,}276\,0730
\end{array}
\]
\[
\begin{aligned}
c&=1^\circ\,26'\,50''{,}3465\qquad
Q\,[&=1^\circ\,4'\,54''{,}4540\,]\\
&&=120199{,}197\,\mathrm{m}.
\end{aligned}
\]
\noindent{\small IX.}\hfill{\small24}\hspace*{.12\linewidth}\par

% Printed186; principalPDF193.
\par\clearpage
\noindent186\hfill{\small\textbf{POSTHUMOUS PAPERS.}}\hfill\null\par\medskip
\[
\begin{array}{@{}l@{\;}r@{\qquad}l@{\;}r@{}}
[Q^3\,.\,.\,.\,.&15{,}239\,705&Q^5\,.\,.\,.\,.&25{,}3995\\
6rr\,.\,.\,.\,.&14{,}385\,912&24r^4\,.\,.&28{,}5957\\\cline{2-2}\cline{4-4}
&0{,}853\,793&&6{,}8038\\
&7{,}1416&&0{,}0006
\end{array}
\]
\[y'-Q=7{,}142\,\mathrm{m}.\]
To obtain \(y'\) from \(Q\), one can also use the following table, Art.~30. Look up the value of \(Q\) in the column for \(x\) and read off the corresponding \(y\).]
\[y'=120206{,}339\,\mathrm{m}.\]
[Further, we have]
\[
\begin{array}{@{}r@{\;}r@{\quad}l@{}}
\operatorname{tang}\tfrac12c\,.\,.\,.&8{,}101\,4346&\\
\tfrac12(Q+y')\,.\,.&5{,}079\,9145&\\\cline{2-2}
&3{,}181\,3491\,.\,.\,.\,.\,.\,.\,.\,.&1518{,}270\,\mathrm{m}\\
&&\Omega=5913075{,}164\\\cline{3-3}
&&x'=5914593{,}434\,\mathrm{m}\,[=53^\circ\,13'\,52''{,}8273].
\end{array}
\]
[If \(x,y\) are the plane coordinates corresponding to the point \(\Phi,\lambda\) on the spheroid, then
\[
\mathrm{f}(x+iy)=\int\frac{(1-ee)\mathrm{d}\Phi}{(1-ee\sin\Phi^2)\cos\Phi}+i\lambda
=\mathrm{f}(\psi)+i\lambda,\quad\text{or}\quad x+iy=\mathrm{g}(\mathrm{f}(\psi)+i\lambda),
\]
\noindent where \(\psi\) denotes the length of the meridian arc from the equator to latitude \(\Phi\). For \(\lambda=0\), we have \(y=0\) and \(x=\psi\). Let \(\omega\) further denote the length of the meridian arc from the equator to latitude \(\Omega\) on the sphere. For \(\lambda=0\), we have \(y'=0\) and \(x'=\omega\). But since
\[
\mathrm{g}(\mathrm{f}(\psi))=\psi=\omega+A\sin2\omega+B\sin4\omega+C\sin6\omega+\ldots
\]
\noindent where, according to p.~179,
\[
A=2ff+\frac{16}{3}f^4+\frac{28}{3}f^6-\ldots,\quad
B=\frac{13}{3}f^4+\frac{464}{15}f^6+\ldots,\quad
C=\frac{244}{15}f^6+\ldots,\text{ and so on,}
\]
\noindent it follows in general that
\[
\mathrm{g}(\mathrm{f}(\psi)+i\lambda)=x+iy=x'+iy'+A\sin2(x'+iy')+B\sin4(x'+iy')+C\sin6(x'+iy')\ldots
\]
If we set
\[\mathfrak{x}=x'+A\sin2x'+B\sin4x'+C\sin6x'+\ldots,\]

% Printed187; principalPDF194.
\par\clearpage
\noindent\null\hfill{\small\textbf{CONFORMAL MAPPING OF THE SPHEROID IN THE PLANE.}}\hfill187\par\medskip
\noindent where \(\mathfrak{x}\) is obtained from the preceding table, p.~184 (for the value of \(x'\) in the column for \(\omega\), look up the corresponding \(\mathfrak{x}\) in the column for \(\psi\)), then
\[
\begin{aligned}
x-\mathfrak{x}+i(y-y')={}&-2A\sin2x'\sin iy^{\prime 2}+iA\cos2x'\frac{\sin2iy'}{i}\\
&-2B\sin4x'\sin2iy^{\prime 2}+iB\cos4x'\frac{\sin4iy'}{i}\\
&-2C\sin6x'\sin3iy^{\prime 2}+iC\cos6x'\frac{\sin6iy'}{i}
\end{aligned}
\]
{\centering etc.\par}
One can also replace \(y'\) by \(Q\). We have
\[\operatorname{tang}Q=\frac{\sin iy'}{i}.]\]
If we set \(\operatorname{tang}Q=t\), then
\[
\begin{aligned}
x={}&\mathfrak{x}+2A\sin2x'.tt\\
&+8B\sin4x'.(tt+t^4)
\end{aligned}
\]
{\centering etc.\par}
\[
\begin{aligned}
y={}&y'+2A\cos2x'.t\sqrt{1+tt}\\
&+4B\cos4x'.(t+2t^3)\sqrt{1+tt}
\end{aligned}
\]
{\centering etc.\par}
[According to p.~183, for the flattening \(\frac1{302{,}7827\ldots}\) we have \(A=5253{,}498\,1263\,\mathrm{m}\), \(B=4{,}700\,0528\,\mathrm{m}\), \(C=0{,}007\,2875\,\mathrm{m}\).

The calculation of \(x,y\) is now to be carried out using the formula given first. For this we have
\[
\log\frac{\sin i.ny'}{i}=\log ny'+\frac{k}{6}(ny')^2-\frac{k}{180}(ny')^4\ldots,\quad n=1,2,3,\ldots,
\]
\noindent where \(y'\) has already been divided by the radius \(r\) of the sphere. Here \(k\) denotes the modulus of the Briggsian logarithms. The calculation of \(\log\frac{\sin i.ny'}{i}\) can also be carried out according to the procedure in Art.~30.]
\[
\begin{array}{@{}l@{\;}r@{\qquad}l@{\;}r@{}}
A\,.\,.\,.\,.\,.\,.\,.&3{,}720\,4486&&\\
\sin2x'\,.\,.\,.&9{,}981\,8206&&\\
\cos2x'\,.\,.\,.&9{,}452\,3858_n&&\\\cline{2-2}
A\sin2x'\,.\,.\,.&3{,}702\,2692&A\cos2x'\,.\,.\,.&3{,}172\,8344_n\\
\sin iy^{\prime 2}\,.\,.\,.\,.&6{,}552\,1460_n&\dfrac{\sin2iy'}{i}\,.\,.\,.\,.&8{,}577\,1804\\\cline{4-4}
2\,.\,.\,.\,.\,.\,.\,.\,.&0{,}301\,0300&&1{,}750\,0148_n\\\cline{2-2}
&0{,}555\,4452&&
\end{array}
\]
\noindent\null\hfill{\small24*}\hspace*{.12\linewidth}\par

% Printed188; principalPDF195.
\par\clearpage
\noindent188\hfill{\small\textbf{POSTHUMOUS PAPERS.}}\hfill\null\par\medskip
\[
\begin{array}{@{}l@{\;}r@{\qquad}l@{\;}r@{}}
B\,.\,.\,.\,.\,.\,.\,.&0{,}672\,1028&&\\
\sin4x'\,.\,.\,.&9{,}735\,2364_n&&\\
\cos4x'\,.\,.\,.&9{,}923\,9583_n&&\\\cline{2-2}
B\sin4x'\,.\,.\,.&0{,}407\,3392_n&B\cos4x'\,.\,.\,.&0{,}596\,0611_n\\
\sin2iy^{\prime 2}\,.\,.\,.&7{,}154\,3608_n&\dfrac{\sin4iy'}{i}\,.\,.\,.\,.&8{,}878\,5200\\\cline{4-4}
2\,.\,.\,.\,.\,.\,.\,.\,.&0{,}301\,0300&&9{,}474\,5811_n\\\cline{2-2}
&7{,}862\,7300&&\\[3pt]
C\,.\,.\,.\,.\,.\,.\,.&7{,}86258&&\\
\sin6x'\,.\,.\,.&9{,}81354_n&&\\
\cos6x'\,.\,.\,.&9{,}88032&&\\\cline{2-2}
C\sin6x'\,.\,.\,.&7{,}67612_n&C\cos6x'\,.\,.\,.&7{,}74290\\
\sin3iy^{\prime 2}\,.\,.\,.&7{,}50680_n&\dfrac{\sin6iy'}{i}\,.\,.\,.\,.&9{,}05513\\\cline{4-4}
2\,.\,.\,.\,.\,.\,.\,.\,.&0{,}30103&&6{,}79803\\\cline{2-2}
&5{,}48395&&
\end{array}
\]
[For]
\[x'=5914593{,}434\,\mathrm{m}\]
\noindent [the corresponding \(\mathfrak{x}\), as already mentioned, is obtained from the table on p.~184 or by using the series for \(\psi-\omega\) in Art.~28.]
\[
\begin{array}{@{}r@{\;}r@{\qquad}r@{\;}r@{}}
\mathfrak{x}=&5919629{,}003\,\mathrm{m}&y'=&120206{,}339\,\mathrm{m}\\
+&3{,}593&-&56{,}236\\
-&0{,}007&-&0{,}298\\
-&0{,}000&+&0{,}001\\\cline{2-2}\cline{4-4}
x=&5919632{,}589\,\mathrm{m}&y=&120149{,}806\,\mathrm{m}\\
\text{Göttingen }[51^\circ\,31'\,48''{,}00]={}&5710161{,}658&&\\\cline{2-2}
\text{Varel}={}&209470{,}931\,\mathrm{m}.&&
\end{array}
\]
\bigskip
{\centering [30.]\par}
The [following] table is used to calculate the hyperbolic functions. If the entries of the two columns are denoted by \(x,y\), then
\[i\operatorname{tang}x=\sin iy\]
\[
\log\operatorname{tang}\tfrac12(90^\circ+x)=\tfrac12k\pi y:10^7
\quad[k=\text{modulus of the Briggsian logarithms}]
\]
\[\log\tfrac12k\pi:10^7=2{,}833\,9042-10.\]

% Printed189; principalPDF196.
\par\clearpage
\noindent\null\hfill{\small\textbf{CONFORMAL MAPPING OF THE SPHEROID IN THE PLANE.}}\hfill189\par\medskip
For larger \(y\) outside the range of the table, one can set:
\[\log\sec\psi=\tfrac12k\pi y:10^7\]
\[\frac{\sin iy}{i}=\frac{\sin\psi^2}{2\cos\psi}=\tfrac12\operatorname{tang}\psi\sin\psi,\]
\noindent or one looks up in my table of logarithms
\[
\begin{array}{@{}r@{\quad}|@{\quad}l@{}}
\begin{aligned}
k\pi y:10^7&=C\\[3pt]
\log\frac{\sin iy}{i}&=\tfrac12C-B-\log2\\
&=A-\tfrac12C-\log2\\
&=\tfrac12(A-B)-\log2
\end{aligned}
&\begin{aligned}
\mathopen{[}k\pi y:10^7={}]&\,B\\[3pt]
[\log\frac{\sin iy}{i}&=-A-\tfrac12B-\log2\\
&=\tfrac12B-C-\log2\\
&=]-\tfrac12(A+C)-\log2
\end{aligned}
\end{array}
\]
We have
\[
\begin{aligned}
y&=x+\tfrac16x^3+\tfrac1{24}x^5+\tfrac{61}{5040}x^7\ldots\\
x&=y-\tfrac16y^3+\tfrac1{24}y^5-\tfrac{61}{5040}y^7\ldots
\end{aligned}
\]
[To calculate the following table, replace \(x\) and \(y\) in these two series by \(\frac{x}{r}\) and \(\frac{y}{r}\), with \(r=\frac{20\,000\,000}{\pi}\).]
{\small\renewcommand{\arraystretch}{.60}
\[\begin{array}{|r|r|r|r|r|}\hline
\rule{0pt}{2.1ex}x&y&\multicolumn{3}{c|}{\text{diff.}}\\
\rule[-3.6pt]{0pt}{3.6pt}\text{(metres)}&\text{(metres)}&\multicolumn{3}{c|}{}\\\hline\hline
\rule{0pt}{2.5ex}0&0&&0&\\
&&10000{,}00411&&2468\\
10\,000&10000{,}00411&&2468&\\
&&02879&&2466\\
20\,000&20000{,}03290&&4934&\\
&&07813&&2469\\
30\,000&30000{,}11103&&7403&\\
&&15216&&2467\\
40\,000&40000{,}26319&&9870&\\
&&25086&&2467\\
50\,000&50000{,}51405&&12337&\\
&&37423&&2469\\
60\,000&60000{,}88828&&14806&\\
&&52229&&2468\\
70\,000&70001{,}41057&&17274&\\
&&69503&&2467\\
80\,000&80002{,}10560&&19741&\\
&&10000{,}89244&&2470\\
90\,000&90002{,}99804&&22211&\\
&&10001{,}11455&&2467\\
100\,000&100004{,}11259&&24678&\\
\cline{1-2}
&&1{,}36133&&2471\\
110\,000&110005{,}47392&&27149&\\
&&1{,}63282&&2469\\
120\,000&120007{,}10674&&29618&\\
&&1{,}92900&&2469\\
130\,000&130009{,}03574&&32087&\\
&&2{,}24987&&2470\\
140\,000&140011{,}28561&&34557&\\
&&2{,}59544&&2472\\
150\,000&150013{,}88105&&37029&\\
&&2{,}96573&&2470\\
160\,000&160016{,}84678&&39499&\\
&&3{,}36072&&2472\\
170\,000&170020{,}20750&&41971&\\
&&3{,}78043&&2472\\
180\,000&180023{,}98793&&44443&\\
&&4{,}22486&&2472\\
190\,000&190028{,}21279&&46915&\\
&&10004{,}69401&&2474\\
200\,000&200032{,}90680&&49389&\\
\end{array}\]
}


% Printed190; principalPDF197.
\par\clearpage
\noindent190\hfill{\small\textbf{POSTHUMOUS PAPERS.}}\hfill\null\par\medskip
{\footnotesize\renewcommand{\arraystretch}{.63}
\[\begin{array}{|r|r|r|r|r|}\hline
\rule{0pt}{2ex}x&y&\multicolumn{3}{c|}{\text{diff.}}\\
\rule[-3.6pt]{0pt}{3.6pt}\text{(metres)}&\text{(metres)}&\multicolumn{3}{c|}{}\\\hline\hline
\rule{0pt}{2.5ex}&&10004{,}69401&&2474\\
200\,000&200032{,}90680&&49389&\\
\cline{1-2}
&&5{,}18790&&2473\\
210\,000&210038{,}09470&&51862&\\
&&5{,}70652&&2475\\
220\,000&220043{,}80122&&54337&\\
&&6{,}24989&&2476\\
230\,000&230050{,}05111&&56813&\\
&&6{,}81802&&2474\\
240\,000&240056{,}86913&&59287&\\
&&7{,}41089&&2478\\
250\,000&250064{,}28002&&61765&\\
&&8{,}02854&&2477\\
260\,000&260072{,}30856&&64242&\\
&&8{,}67096&&2477\\
270\,000&270080{,}97952&&66719&\\
&&10009{,}33815&&2480\\
280\,000&280090{,}31767&&69199&\\
&&10010{,}03014&&2480\\
290\,000&290100{,}34781&&71679&\\
&&10{,}74693&&2480\\
300\,000&300111{,}09474&&74159&\\
\cline{1-2}
&&11{,}48852&&2481\\
310\,000&310122{,}58326&&76640&\\
&&12{,}25492&&2484\\
320\,000&320134{,}83818&&79124&\\
&&13{,}04616&&2484\\
330\,000&330147{,}88434&&81608&\\
&&13{,}86224&&2483\\
340\,000&340161{,}74658&&84091&\\
&&14{,}70315&&2485\\
350\,000&350176{,}44973&&86576&\\
&&15{,}56891&&2488\\
360\,000&360192{,}01864&&89064&\\
&&16{,}45955&&2488\\
370\,000&370208{,}47819&&91552&\\
&&17{,}37507&&2488\\
380\,000&380225{,}85326&&94040&\\
&&18{,}31547&&2491\\
390\,000&390244{,}16873&&96531&\\
&&19{,}28078&&2490\\
400\,000&400263{,}44951&&99021&\\
\cline{1-2}
&&20{,}27099&&2493\\
410\,000&410283{,}72050&&101514&\\
&&21{,}28613&&2494\\
420\,000&420305{,}00663&&104008&\\
&&22{,}32621&&2495\\
430\,000&430327{,}33284&&106503&\\
&&23{,}39124&&2496\\
440\,000&440350{,}72408&&108999&\\
&&24{,}48123&&2498\\
450\,000&450375{,}20531&&111497&\\
&&25{,}59620&&2499\\
460\,000&460400{,}80151&&113996&\\
&&26{,}73616&&2501\\
470\,000&470427{,}53767&&116497&\\
&&27{,}90113&&2501\\
480\,000&480455{,}43880&&118998&\\
&&29{,}09111&&2504\\
490\,000&490484{,}52991&&121502&\\
&&30{,}30613&&2504\\
500\,000&500514{,}83604&&124006&\\
\cline{1-2}
&&31{,}54619&&2508\\
510\,000&510546{,}38223&&126514&\\
&&32{,}81133&&2507\\
520\,000&520579{,}19356&&129021&\\
&&34{,}10154&&2510\\
530\,000&530613{,}29510&&131531&\\
&&35{,}41685&&2511\\
540\,000&540648{,}71195&&134042&\\
&&36{,}75727&&2513\\
550\,000&550685{,}46922&&136555&\\
&&38{,}12282&&2515\\
560\,000&560723{,}59204&&139070&\\
&&39{,}51352&&2516\\
570\,000&570763{,}10556&&141586&\\
&&40{,}92938&&2518\\
580\,000&580804{,}03494&&144104&\\
&&42{,}37042&&2521\\
590\,000&590846{,}40536&&146625&\\
&&10043{,}83667&&2521\\
600\,000&600890{,}24203&&149146&\\
\end{array}\]
}


% Printed191; principalPDF198.
\par\clearpage
\noindent\null\hfill{\small\textbf{CONFORMAL MAPPING OF THE SPHEROID IN THE PLANE.}}\hfill191\par\medskip
{\small\renewcommand{\arraystretch}{.85}
\[\begin{array}{|r|r|r|r|r|}\hline
x&y&\multicolumn{3}{c|}{\text{diff.}}\\
\text{(metres)}&\text{(metres)}&\multicolumn{3}{c|}{}\\\hline\hline
&&10043{,}83667&&2521\\
600\,000&600890{,}24203&&149146&\\
\cline{1-2}
&&45{,}32813&&2524\\
610\,000&610935{,}57016&&151670&\\
&&46{,}84483&&2525\\
620\,000&620982{,}41499&&154195&\\
&&48{,}38678&&2528\\
630\,000&631030{,}80177&&156723&\\
&&10049{,}95401&&\\
640\,000&641080{,}75578&&&\\
\end{array}\]
}

\bigskip
{\centering [31.]\par\medskip
[Calculation of longitude and latitude from the plane coordinates.]\par}
If
\[p=q-h\sin q,\]
\noindent then [for hyperbolic logarithms] we have
\[
\begin{aligned}
\log\operatorname{tang}\tfrac12p={}&\log\operatorname{tang}\tfrac12q-h-\tfrac12hh\cos q\\
&-\tfrac16h^3(2-\sin q^2)-\tfrac1{24}h^4\cos q(6-\sin q^2)\ldots.
\end{aligned}
\]
If
\[h=a\cos q+b\cos3q+c\cos5q+d\cos7q+\ldots,\]
\noindent then
\[
\begin{aligned}
\log\operatorname{tang}\tfrac12p={}&\log\operatorname{tang}\tfrac12q\\
&-\cos q\left(a+\tfrac38aa+\tfrac14ab+\tfrac14bb+\tfrac{11}{48}a^3+\tfrac9{32}aab+\tfrac{235}{1536}a^4\ldots\right)\\
&-\cos3q\left(b+\tfrac18aa+\tfrac12ab+\tfrac14ac+\tfrac3{32}a^3+\tfrac7{16}aab+\tfrac{121}{1536}a^4\ldots\right)\\
&-\cos5q\left(c+\tfrac14ab+\tfrac18bb+\tfrac12ac+\tfrac1{96}a^3+\tfrac14aab+\tfrac9{512}a^4\ldots\right)\\
&-\cos7q\left(d+\tfrac18bb+\tfrac14ac+\tfrac1{32}aab+\tfrac1{1536}a^4\ldots\right)
\end{aligned}
\]
{\centering etc.\par}
[Set \(p=90^\circ+\omega\) and \(q=90^\circ+\psi\). Starting from the last series relating \(\omega\) and \(\psi\) in Art.~25, one first obtains:]
\[
\begin{aligned}
a&=-4ff-10f^4-\frac{40}{3}f^6+\frac{1247}{9}f^8\ldots\\
b&=\phantom{-4ff-}\frac23f^4+\frac{44}{3}f^6+\frac{1897}{15}f^8\ldots\\
c&=\phantom{-4ff-10f^4}-\frac{68}{15}f^6-\frac{3253}{45}f^8\ldots\\
d&=\phantom{-4ff-10f^4-\frac{40}{3}f^6}+\frac{4397}{315}f^8\ldots
\end{aligned}
\]
{\centering etc.\par}
